\documentclass[11pt,a4paper]{article}
\usepackage[UTF8]{ctex}
\usepackage{geometry}
\geometry{left=2.5cm,right=2.5cm,top=2.5cm,bottom=2.5cm}
\usepackage{amsmath,amssymb,amsthm}
\usepackage{hyperref}
\usepackage{booktabs}
\usepackage{enumitem}
\usepackage{xltabular}
\hypersetup{colorlinks=true,linkcolor=blue,urlcolor=blue}

\newtheorem{definition}{定义}[section]
\newtheorem{proposition}{命题}[section]
\newtheorem{remark}{注记}[section]

\title{多维联合优化收敛性的深化研究（修正版）\\
——联合凸性、退化效应与离散粒度的贪心分析}
\author{李龙胜}
\date{\today}

\begin{document}

\maketitle

\begin{center}
\textit{
贪心择优在单维度上已被严格证明等价于全局最优。\\
但当质量与速率两个维度共享同一块预算时，\\
贪心交替分配能否收敛到全局联合最优解？\\
答案依赖于一个未显式验证的前提——联合凸性。\\
本文精确刻画这个证明缺口，\\
将退化效应与离散粒度纳入统一的联合收敛性框架，\\
并对外部审查中识别的耦合函数形式、鞍点条件和误差紧致性
逐条补充了参数化分析和边界讨论。
}
\end{center}

\vspace{5mm}

\begin{abstract}
贪心择优定理在AI博弈维度的严格证明（单独归档）
已完成了质量维度和速率维度的单维贪心最优性证明。
但在多维联合优化中，贪心交替分配预算的过程
是否收敛到全局最优解，依赖于一个未显式验证的前提——
各维度的成本函数加性可分，或联合成本函数满足联合凸性。
本文对该前提进行深化分析：
建立联合成本函数的耦合结构，
讨论联合凸性在大参数范围内成立但在极端配置下可能出现鞍点；
将退化效应作为负增量纳入联合框架，提出贪心规避策略；
分析离散粒度对交换论证的修正，给出联合误差的近似上界。
本文经外部审查修正，补充了速率压制指数衰减的物理下限、
鞍点条件的Hessian参数化方向以及离散粒度误差上界的紧致性分析。
\end{abstract}

\tableofcontents

\section{引言：从单维最优到联合最优的鸿沟}

贪心择优定理在AI博弈维度的严格证明
已在独立文档中完成。
质量维度和速率维度的单维贪心最优性均已获得严格证明。

然而，在真实AI化攻防中，
质量维度和速率维度共享同一块分析能力预算 \(C\)。
防御方需要在两个维度之间交替分配有限资源。

联合优化的收敛性要求证明：
贪心在两个维度之间交替分配预算的过程，
收敛到的解确实是全局联合最优解，
而不是两个独立最优解的简单叠加。
这个证明依赖于一个未显式验证的前提——
各维度的成本函数加性可分，
或联合成本函数满足联合凸性。
本文对该前提进行深化分析。

\section{联合成本函数的耦合结构}

设防御方在质量维度的投入为 \(x\)，
在速率维度的投入为 \(y\)，
总预算约束为 \(x + y \le C\)。

联合成本函数为
\[
J(x, y) = J_{\text{quality}}(x) + J_{\text{rate}}(y) + J_{\text{coupling}}(x, y),
\]
其中 \(J_{\text{quality}}(x)\) 是质量维度的独立成本，
\(J_{\text{rate}}(y)\) 是速率维度的独立成本，
\(J_{\text{coupling}}(x, y)\) 是两个维度的耦合成本。

\subsection{耦合成本的来源}

速率投入增加压缩了攻击方探测速率 \(R_A\)，
降低了防御方在质量维度上所需的对抗训练频率。
两个维度的边际成本通过攻击方的行为间接耦合。

设攻击方探测速率 \(R_A\) 受防御方速率压制强度 \(y\) 的影响：
\(R_A = R_A^0 \cdot e^{-\gamma y}\)，其中 \(\gamma > 0\) 是压制系数。
防御方在质量维度上所需的对抗训练频率
与攻击方探测速率成正比——攻击方探测越快，防御方需训练越频繁。
因此，质量维度的边际成本随速率投入的增加而下降，
耦合成本 \(J_{\text{coupling}}(x, y)\) 为负值。

\textbf{指数衰减模型的物理下限}：
当前模型 \(R_A = R_A^0 \cdot e^{-\gamma y}\) 在 \(y \to \infty\) 时趋于零。
在真实对抗中，攻击方探测速率不可能被压到零——
攻击方可通过更换IP、使用代理、降低探测频率来对抗压制，
但始终维持一个非零下限 \(R_A^{\min}\)。
在 \(y\) 较小时，指数衰减是合理的近似；
在 \(y\) 较大时，模型需要修正以反映下限的存在。
当前指数衰减模型在参数校准后可作为一阶近似。

\subsection{加性可分与联合凸性}

若耦合成本为零，即 \(J_{\text{coupling}}(x, y) = 0\)，
联合成本函数加性可分，
联合最优解退化为两个独立最优解的叠加，
贪心交替分配自然收敛到全局最优。

若耦合成本不为零，
联合凸性要求 \(J(x, y)\) 关于 \((x, y)\) 是联合凸函数。
在大参数范围内，\(J_{\text{quality}}(x)\) 和 \(J_{\text{rate}}(y)\)
各自的严格凸性保证了联合凸性成立——
严格凸函数的和仍然是严格凸函数，
耦合成本在大多数区域内是线性的或凹的，
不会破坏联合凸性。

但在某些极端配置下可能出现鞍点：
速率投入极高且质量投入极低时，
继续增加速率投入的边际收益趋近于零
（因为攻击方探测速率已被压到物理下限），
而质量投入的边际收益仍然很高，
联合成本函数在此处可能出现局部鞍点——
两个维度的替代弹性趋近于零，
联合凸性被弱化。

\textbf{鞍点条件的参数化方向}：
若联合成本函数在局部满足 \(\det(H) \le 0\)
——其中 \(H\) 是 \(J(x, y)\) 的Hessian矩阵——
联合凸性在该点局部失效。
将具体的耦合成本函数代入该判别式，
可解析求解鞍点区域的参数边界。
当前版本给出了鞍点的定性条件，
定量边界待成本函数实测标定后计算。

\section{退化效应在联合框架中的位置}

在AI化攻防的真实场景中，
防御方的对抗训练可能导致正常样本精度下降，
攻击方的高质量对抗样本可能因防御方特征空间微调而过期。
这些退化效应在联合成本函数中表现为
质量进展的负增量——
某些训练步骤可能产生负的边际收益。

设防御方在第 \(t\) 步训练后，
正常样本上的检测精度下降 \(\Delta q_D^{\text{deg}} > 0\)。
这一退化效应使质量维度的实际进展量
减少 \(\Delta q_D^{\text{deg}}\)，
在成本函数中体现为额外的退化成本 \(C_{\text{deg}}(\Delta q_D^{\text{deg}})\)。

退化效应的存在不破坏联合凸性在大参数范围内的成立，
但在退化密集区域可能增加局部鞍点的风险——
质量维度的边际收益在此处被退化成本抵消，
贪心选择可能倾向于回避质量维度，
过度向速率维度分配预算。

防御方需要在贪心选择中引入对退化的预期：
当退化风险超过某阈值时，
跳过该步骤的贪心选择，
转入速率维度或其他维度进行补偿。
这一策略在贪心框架中自然成立——
每一步贪心选择的对象
从“当前维度的最佳操作”
扩展为“所有可行维度中的最佳操作”。

\section{离散粒度对交换论证的修正}

在严格证明中，交换论证假设进展量可以连续分配——
均匀分配总进展量使总成本最低。
但在真实场景中，训练操作是离散的——
每次对抗训练要么执行，要么不执行。
进展量无法均匀分配。

设每次训练操作提升的检测质量为固定的离散粒度 \(\Delta q_D^{\text{unit}}\)。
总目标提升量 \(Q_D\) 需要 \(n = \lceil Q_D / \Delta q_D^{\text{unit}} \rceil\) 次操作完成。
贪心路径与全局最优路径之间的成本差
受离散化粒度的约束：
\[
|J_D(\pi_{\text{greedy}}) - J_D(\pi_{\text{opt}})| \le
C_D'' \cdot (\Delta q_D^{\text{unit}})^2 \cdot n,
\]
其中 \(C_D''\) 是成本函数二阶导数的上界。

当离散粒度相对于总进展量足够小时，
成本差趋于零，贪心路径近似最优。

\textbf{误差上界的紧致性}：
当前上界使用全局 \(C_D''\) 的最大值，可能宽松——
实际贪心路径和最优路径之间的成本差可能远小于此上界。
更紧的上界可通过以下两种路径获得：
使用 \(C_D''\) 在贪心路径上的局部值替代全局最大值；
或通过分段线性逼近收紧误差估计。
更紧上界的推导待进一步数学工作。
当前宽松上界不影响近似最优性的定性结论。

在多维联合优化中，
两个维度各自的离散化误差可能叠加。
联合误差的绝对上界是两个维度独立误差上界之和。
当两个维度的离散粒度都足够小时，
贪心路径在联合空间中同样近似最优。

\section{诚实标注}

\begin{center}
\begin{xltabular}{\textwidth}{c|X|c}
\toprule
\textbf{命题} & \textbf{状态} & \textbf{层级} \\
\midrule
联合收敛性的严格证明 &
联合凸性在大部分参数范围内定性成立，
严格验证需成本函数的实测数据标定 &
II（联合凸性的精确参数区域待标定） \\
速率压制指数衰减的物理下限 &
指数衰减在 \(y\) 较小时合理，
\(y\) 较大时需修正以反映非零下限 &
II（修正函数形式待实测数据拟合） \\
联合凸性的鞍点风险 &
在极端配置下可能出现局部鞍点，
已给出Hessian判别式的参数化方向 &
II（定量边界待成本函数具体形式确定后计算） \\
退化效应在联合框架中的建模 &
已在联合成本函数中引入负增量，
贪心规避策略已提出 &
II（退化成本的具体函数形式需实证数据） \\
离散粒度对交换论证的修正 &
近似误差上界已给出，
贪心路径在离散空间中近似最优 &
II（更紧上界待进一步数学工作） \\
多维联合优化的贪心最优性 &
单维度严格证明已完成，
联合收敛性依赖联合凸性条件 &
II（条件验证和严格收敛性证明待后续工作） \\
\bottomrule
\end{xltabular}
\end{center}

\section{结语}

多维联合优化收敛性的深化研究已经完成，
并经外部审查修正。
联合收敛性的核心前提——联合成本函数的联合凸性——
在大参数范围内定性成立，但在极端配置下可能出现鞍点。
速率压制的指数衰减模型在物理下限附近需要修正，
鞍点条件可通过Hessian判别式参数化，
离散粒度误差上界可通过局部二阶导数值进一步收紧。
退化效应作为负增量被纳入联合框架，
通过贪心规避策略得到处理。
贪心择优定理在AI博弈维度的严格证明
因此获得了对其边界条件的精确刻画。

\vspace{5mm}
\noindent\textbf{文档信息}：
\begin{itemize}
    \item 作者：李龙胜
    \item 版本：修正版（整合外部审查反馈）
    \item 许可：CC BY 4.0
    \item 前置依赖：贪心择优定理在AI博弈维度的严格证明、
          弹药质量维度深化修正版、
          柯克霍夫边界条件与速率博弈修正版
\end{itemize}

\end{document}